%%
%% This is file `sample-sigconf.tex',
%% generated with the docstrip utility.
%%
%% The original source files were:
%%
%% samples.dtx  (with options: `all,proceedings,bibtex,sigconf')
%% 
%% IMPORTANT NOTICE:
%% 
%% For the copyright see the source file.
%% 
%% Any modified versions of this file must be renamed
%% with new filenames distinct from sample-sigconf.tex.
%% 
%% For distribution of the original source see the terms
%% for copying and modification in the file samples.dtx.
%% 
%% This generated file may be distributed as long as the
%% original source files, as listed above, are part of the
%% same distribution. (The sources need not necessarily be
%% in the same archive or directory.)
%%
%%
%% Commands for TeXCount
%TC:macro \cite [option:text,text]
%TC:macro \citep [option:text,text]
%TC:macro \citet [option:text,text]
%TC:envir table 0 1
%TC:envir table* 0 1
%TC:envir tabular [ignore] word
%TC:envir displaymath 0 word
%TC:envir math 0 word
%TC:envir comment 0 0
%%
%% The first command in your LaTeX source must be the \documentclass
%% command.
%%
%% For submission and review of your manuscript please change the
%% command to \documentclass[manuscript, screen, review]{acmart}.
%%
%% When submitting camera ready or to TAPS, please change the command
%% to \documentclass[sigconf]{acmart} or whichever template is required
%% for your publication.
%%
%%
% \documentclass[sigconf]{acmart}
\documentclass[sigconf,natbib=true]{acmart}
\usepackage{enumitem}
\usepackage{multirow}
\usepackage{booktabs}
\usepackage{subcaption} % support subfigure env
\usepackage{hyperref}
\usepackage{colortbl}
\usepackage{graphicx}   % 插入图片
\usepackage{makecell}   % 支持单元格多行
\usepackage{xcolor}


%%
%% \BibTeX command to typeset BibTeX logo in the docs
\AtBeginDocument{%
  \providecommand\BibTeX{{%
    Bib\TeX}}}

%% Rights management information.  This information is sent to you
%% when you complete the rights form.  These commands have SAMPLE
%% values in them; it is your responsibility as an author to replace
%% the commands and values with those provided to you when you
%% complete the rights form.
% \setcopyright{acmlicensed}
% \copyrightyear{2026}
% \acmYear{2026}
% \acmDOI{XXXXXXX.XXXXXXX}
%% These commands are for a PROCEEDINGS abstract or paper.
% \acmConference[SIGIR '26]{the 49th International ACM SIGIR Conference on Research and Development in Information Retrieval}{July 20--24,
%   2026}{Melbourne, Australia}
%%
%%  Uncomment \acmBooktitle if the title of the proceedings is different
%%  from ``Proceedings of ...''!
%%
%%\acmBooktitle{Woodstock '18: ACM Symposium on Neural Gaze Detection,
%%  June 03--05, 2018, Woodstock, NY}
% \acmISBN{978-1-4503-XXXX-X/2018/06}

\copyrightyear{2026}
\acmYear{2026}
\setcopyright{cc}
\setcctype{by}
\acmConference[SIGIR '26] {Proceedings of the 49th International ACM SIGIR Conference on Research and Development in Information Retrieval}{July 20--24, 2026}{Melbourne, VIC, Australia.}
\acmBooktitle{Proceedings of the 49th International ACM SIGIR Conference on Research and Development in Information Retrieval (SIGIR '26), July 20--24, 2026, Melbourne, VIC, Australia}
\acmISBN{979-8-4007-2599-9/2026/07}
\acmDOI{10.1145/3805712.3809892}
% 1 Authors, replace the red X's with your assigned DOI string during the rightsreview eform process.
% 2 Your DOI link will become active when the proceedings appears in the DL.
% 3 Retain the DOI string between the curly braces for uploading your promo video.

\settopmatter{printacmref=true}


%%
%% Submission ID.
%% Use this when submitting an article to a sponsored event. You'll
%% receive a unique submission ID from the organizers
%% of the event, and this ID should be used as the parameter to this command.
%%\acmSubmissionID{123-A56-BU3}

%%
%% For managing citations, it is recommended to use bibliography
%% files in BibTeX format.
%%
%% You can then either use BibTeX with the ACM-Reference-Format style,
%% or BibLaTeX with the acmnumeric or acmauthoryear sytles, that include
%% support for advanced citation of software artefact from the
%% biblatex-software package, also separately available on CTAN.
%%
%% Look at the sample-*-biblatex.tex files for templates showcasing
%% the biblatex styles.
%%

%%
%% The majority of ACM publications use numbered citations and
%% references.  The command \citestyle{authoryear} switches to the
%% "author year" style.
%%
%% If you are preparing content for an event
%% sponsored by ACM SIGGRAPH, you must use the "author year" style of
%% citations and references.
%% Uncommenting
%% the next command will enable that style.
%%\citestyle{acmauthoryear}


%%
%% end of the preamble, start of the body of the document source.
\begin{document}

%%
%% The "title" command has an optional parameter,
%% allowing the author to define a "short title" to be used in page headers.
\title{RoTE: Coarse-to-Fine Multi-Level Rotary Time Embedding for Sequential Recommendation}

%%
%% The "author" command and its associated commands are used to define
%% the authors and their affiliations.
%% Of note is the shared affiliation of the first two authors, and the
%% "authornote" and "authornotemark" commands
%% used to denote shared contribution to the research.

\author{Haolin Zhang}
\authornote{Both authors contributed equally to this research.}
\affiliation{%
  \institution{Shenzhen International Graduate School, Tsinghua University}
  \city{Shenzhen}
  \country{China}
}
\email{zhanghaolin24@mails.tsinghua.edu.cn}

\author{Longtao Xiao}
\authornotemark[1]
\affiliation{%
  \institution{Huazhong University of Science and Technology}
  \city{Wuhan}
  \country{China}
}
\email{xiaolongtao@hust.edu.cn}

\author{Guohao Cai}
\affiliation{%
  \institution{Huawei Noah’s Ark Lab}
  \city{Shenzhen}
  \country{China}
}
\email{caiguohao1@huawei.com}

\author{Ruixuan	Li}
\affiliation{%
  \institution{Huazhong University of Science and Technology}
  \city{Wuhan}
  \country{China}
}
\email{rxli@hust.edu.cn}



\author{Xiu Li}
\affiliation{%
  \institution{Shenzhen International Graduate School, Tsinghua University}
  \city{Shenzhen}
  \country{China}
}
\email{li.xiu@sz.tsinghua.edu.cn}
\authornote{Corresponding author}

%%
%% By default, the full list of authors will be used in the page
%% headers. Often, this list is too long, and will overlap
%% other information printed in the page headers. This command allows
%% the author to define a more concise list
%% of authors' names for this purpose.
\renewcommand{\shortauthors}{Zhang et al.}

%%
%% The abstract is a short summary of the work to be presented in the
%% article.
\begin{abstract}
Sequential recommendation models have been widely adopted for modeling user behavior. Existing approaches typically construct user interaction sequences by sorting items according to timestamps and then model user preferences from historical behaviors. While effective, such a process only considers the order of temporal information but overlooks the actual time spans between interactions, resulting in a coarse representation of users’ temporal dynamics and limiting the model’s ability to capture long-term and short-term interest evolution. To address this limitation, we propose \textbf{RoTE}, a novel multi-level temporal embedding module that explicitly models time span information in sequential recommendation. RoTE decomposes each interaction timestamp into multiple temporal granularities, ranging from coarse to fine, and incorporates the resulting temporal representations into item embeddings. This design enables models to capture heterogeneous temporal patterns and better perceive temporal distances among user interactions during sequence modeling. RoTE is a lightweight, plug-and-play module that can be seamlessly integrated into existing Transformer-based sequential recommendation models without modifying their backbone architectures. We apply RoTE to several representative models and conduct extensive experiments on three public benchmarks. Experimental results demonstrate that RoTE consistently enhances the corresponding backbone models, achieving up to a 20.11\% improvement in NDCG@5, which confirms the effectiveness and generality of the proposed approach. Our code is available at https://github.com/XiaoLongtaoo/RoTE.

 
\end{abstract}

%%
%% The code below is generated by the tool at http://dl.acm.org/ccs.cfm.
%% Please copy and paste the code instead of the example below.
%%
\begin{CCSXML}
<ccs2012>
 <concept>
  <concept_id>00000000.0000000.0000000</concept_id>
  <concept_desc>Do Not Use This Code, Generate the Correct Terms for Your Paper</concept_desc>
  <concept_significance>500</concept_significance>
 </concept>
 <concept>
  <concept_id>00000000.00000000.00000000</concept_id>
  <concept_desc>Do Not Use This Code, Generate the Correct Terms for Your Paper</concept_desc>
  <concept_significance>300</concept_significance>
 </concept>
 <concept>
  <concept_id>00000000.00000000.00000000</concept_id>
  <concept_desc>Do Not Use This Code, Generate the Correct Terms for Your Paper</concept_desc>
  <concept_significance>100</concept_significance>
 </concept>
 <concept>
  <concept_id>00000000.00000000.00000000</concept_id>
  <concept_desc>Do Not Use This Code, Generate the Correct Terms for Your Paper</concept_desc>
  <concept_significance>100</concept_significance>
 </concept>
</ccs2012>
\end{CCSXML}

\ccsdesc[500]{Information systems~Recommender systems}
% \ccsdesc[300]{Do Not Use This Code~Generate the Correct Terms for Your Paper}
% \ccsdesc{Do Not Use This Code~Generate the Correct Terms for Your Paper}
% \ccsdesc[100]{Do Not Use This Code~Generate the Correct Terms for Your Paper}

%%
%% Keywords. The author(s) should pick words that accurately describe
%% the work being presented. Separate the keywords with commas.
\keywords{Sequential Recommendation, Temporal Dynamics, Plug-and-Play Module}
%% A "teaser" image appears between the author and affiliation
%% information and the body of the document, and typically spans the
%% page.
% \begin{teaserfigure}
%   \includegraphics[width=\textwidth]{sampleteaser}
%   \caption{Seattle Mariners at Spring Training, 2010.}
%   \Description{Enjoying the baseball game from the third-base
%   seats. Ichiro Suzuki preparing to bat.}
%   \label{fig:teaser}
% \end{teaserfigure}

% \received{20 February 2007}
% \received[revised]{12 March 2009}
% \received[accepted]{5 June 2009}

%%
%% This command processes the author and affiliation and title
%% information and builds the first part of the formatted document.
\maketitle

\section{Introduction}

% Sequential recommendation has become a fundamental research problem in recommender systems, aiming to predict users’ future interactions based on their historical behavior sequences. By modeling interaction histories as sorted item sequences, sequential recommender models can effectively capture dynamic user preferences and evolving behavioral patterns. Existing sequential recommendation methods can be broadly categorized into traditional (discriminative) approaches and generative approaches \cite{xiao2025unger}. 

% Sequential recommendation has become a fundamental research problem in recommender systems \cite{pan2026survey, wei2025sequential}, aiming to predict users’ future interactions based on their historical behavior sequences. Sequential recommender models  \cite{fpmc, covington2016deep, shah2025dsrs, he2016fusing} can capture dynamic user preferences by modeling users’ interaction sequences, which are constructed by chronologically ordering the items a user has interacted with. Existing sequential recommendation methods can be broadly categorized into traditional (discriminative) approaches \cite{elecrec} and generative approaches \cite{xiao2025unger}. 


% Traditional methods typically formulate sequential recommendation as a next-item prediction problem  \cite{seqrec_study}, where the model learns to rank candidate items conditioned on the user’s interaction history. Early works adopt recurrent \cite{gru4rec} or convolutional \cite{caser} architectures to model sequential dependencies, while recent approaches predominantly rely on Transformer-based models \cite{sasrec, bert4rec}. More recently, generative sequential recommendation \cite{www25-gen-rec-tutorial} has emerged as an alternative paradigm. Instead of directly ranking items, generative methods model user behavior sequences in an autoregressive or sequence-to-sequence manner, treating recommendation as a generation task \cite{hou2025survey}. By leveraging large language models \cite{p5,gptrec} or sequence generation frameworks \cite{tiger, hstu}, such approaches aim to unify representation learning and prediction, offering greater modeling flexibility and extensibility. 

Sequential recommendation \cite{seqrec_study, pan2026survey, wei2025sequential} is a fundamental problem in recommender systems, which aims to predict the next item a user will interact with based on their historical behavior sequences. Sequential recommender models \cite{fpmc, covington2016deep, shah2025dsrs, he2016fusing} capture dynamic user preferences by modeling interaction sequences, constructed by chronologically ordering previously interacted items. Existing methods can be broadly categorized into traditional (discriminative) approaches \cite{elecrec} and generative approaches \cite{xiao2025unger}. Traditional sequential recommendation \cite{gru4rec, caser, sasrec, bert4rec} formulate the recommendation task as a candidate matching task, in which the model learns user preferences from historical interaction sequences and derives latent user interest representations, which are subsequently utilized to calculate the relevance of candidate items, thus guiding the ranking and selection of items. Generative sequential recommendation \cite{www25-gen-rec-tutorial, hou2025survey} treats recommendation as a sequence generation problem, modeling user behavior in an autoregressive or sequence-to-sequence manner and often leveraging large language models \cite{p5, gptrec} or sequence generation frameworks \cite{tiger, hstu} to achieve more flexible and unified representation learning.

% Despite the strong empirical performance of traditional and generative sequential recommendation models, a key limitation persists in how they handle temporal information. Most existing traditional and generative methods typically sort interaction sequences by timestamp and then rely on positional embeddings or relative order as a coarse proxy for temporal information to capture sequence structure, where such positional encoding can be viewed as a highly simplified form of temporal modeling that implicitly assumes uniformly spaced interactions. This assumption ignores the true time intervals between interactions, results in an oversimplified representation of temporal dynamics and inhibits the model’s ability to capture nuanced temporal patterns such as drift in preferences over irregular time gaps.

Despite the strong empirical performance of traditional and generative sequential recommendation models, a key limitation \cite{nasir2023survey} persists in how they handle temporal information. Most existing traditional and generative methods typically sort interaction sequences by timestamp and then rely on positional embeddings to capture temporal information. Such positional encoding \cite{lopez2024positional} can be viewed as the most simplified form of temporal modeling, as it primarily preserves relative order while treating time as uniformly spaced and ignoring heterogeneous time intervals between interactions. As a result, rich temporal signals at different time scales are either entangled or discarded, leading to an oversimplified representation of temporal dynamics and limiting the model’s ability to capture nuanced patterns such as multi-scale preference evolution and drift over irregular time gaps.

In real-world scenarios, interaction time is commonly recorded in the form of Unix timestamps \cite{hauser2018unix}, providing a precise and continuous measure of when user actions occur. However, such representations do not align well with how users perceive and organize time in daily life. Instead of reasoning in terms of absolute numerical values, users naturally interpret time in a hierarchical and semantic manner \cite{poppel1997hierarchical}, where long-term interests evolve across years, medium-term behaviors are shaped by seasonal or monthly patterns, and short-term intents are driven by daily interactions. This mismatch between the raw timestamp representation and human time perception introduces a bias in temporal modeling, as heterogeneous temporal signals at different scales are compressed into a single continuous value. 

To address these challenges, we propose \textbf{Ro}tary \textbf{T}ime \textbf{E}mbedding (RoTE), a lightweight and plug-and-play multi-level temporal modeling module for sequential recommendation. RoTE explicitly captures coarse-to-fine temporal information and encodes them into item embeddings in the self-attention module. Specifically, RoTE decomposes interaction timestamps into year, month, and day, and integrates them into item representations through a unified rotary embedding mechanism. For each temporal level, RoTE injects time-aware signals directly into the item embedding space, enabling temporal information to be tightly coupled with sequence modeling rather than treated as auxiliary features. The resulting time-enhanced representations are then aggregated to form the final item embeddings used by the model. By seamlessly integrating multi-level temporal signals into the self-attention process without modifying backbone architectures, RoTE allows existing Transformer-based sequential recommendation models to better perceive heterogeneous temporal distances among user interactions without introducing significant computational overhead.

The key contributions of our work are as follows:

\begin{itemize}[left=0pt]

    \item We propose RoTE, a novel plug-and-play temporal embedding module that explicitly models time-span information for sequential recommendation.

    \item We introduce a coarse-to-fine temporal encoding strategy that enables Transformer-based sequential recommendation models to capture multi-level temporal patterns.

    \item We demonstrate that RoTE consistently improves recommendation performance through extensive experiments on multiple benchmark datasets and representative models, showcasing its advantages in performance and generality.
    
\end{itemize}


\section{Methodology}

\subsection{Problem Definition}

Let $\mathcal{U}$ and $\mathcal{I}$ denote the sets of users and items, respectively.
Each user $u \in \mathcal{U}$ is associated with an interaction sequence
\begin{equation}
\mathcal{S}_u = \{(i_0, t_0), (i_1, t_1), \dots, (i_{L_u-1}, t_{L_u-1})\},
\end{equation}
where $i_k \in \mathcal{I}$ denotes the $k$-th interacted item and $t_k$ is the corresponding interaction timestamp, satisfying $t_0 < t_1 < \dots < t_{L_u-1}$.
The sequence length $L_u$ may vary across users, and the time intervals between consecutive interactions are generally irregular.

Given a user’s historical interaction sequence $\mathcal{S}_u$, the goal of sequential recommendation is to model the dependency between past interactions and future behaviors, and predict the next item $i_{L_u}$ conditioned on the observed sequence.
Formally, the task can be defined as learning a conditional distribution
\begin{equation}
p(i_{L_u} \mid \mathcal{S}_u),
\end{equation}
which can be instantiated either as a ranking-based scoring function or as a generative model that directly produces the next item token.

In this work, we focus on enhancing sequential recommendation models by more effectively modeling the temporal information embedded in interaction timestamps, as shown in Figure \ref{fig:RoTE}.
Specifically, we aim to capture heterogeneous time intervals and multi-scale temporal dynamics in a unified manner that is compatible with Transformer-based generative and discriminative sequential recommendation architectures.

\subsection{Temporal Feature Construction}

In real-world scenarios, interaction time is typically recorded as a Unix timestamp. To obtain a structured temporal representation, we decompose each timestamp $t_k$ into a hierarchical triplet
\begin{equation}
t_k \rightarrow (y_k, m_k, d_k),
\end{equation}
where $y_k$, $m_k$, and $d_k$ denote the cumulative number of years, months, and days elapsed since January 1, 1970, respectively.
For example, January 1, 1970 corresponds to $(0, 0, 0)$, while January 1, 1971 corresponds to $(1, 12, 365)$. 
This ordinal representation preserves temporal ordering and enables modeling temporal dynamics across multiple granularities, providing a foundation for for temporal encoding.


\begin{figure}
    \centering
    \includegraphics[width=\linewidth]{SIGIR2026/figures/RoTE_v2.2.pdf}
    \caption{Illustration of the proposed RoTE module.}
    \label{fig:RoTE}
\end{figure}

\subsection{Rotary Time Embedding (RoTE)}

Based on the constructed multi-level temporal features, we propose \textbf{Ro}tary \textbf{T}ime \textbf{E}mbedding (RoTE), a lightweight temporal modeling module that injects time information directly into item representations within the self-attention mechanism. RoTE integrates time-aware signals into the item embedding space by applying rotary-style transformations \cite{rope} to the query and key representations, allowing the attention mechanism to perceive heterogeneous temporal distances among interactions.

Given an input sequence of item embeddings, we first project them into query, key, and value representations following standard multi-head self-attention.
For each interaction, its timestamp is represented by a hierarchical  triplet $(y_k, m_k, d_k)$.
For each temporal level, RoTE computes rotation angles using a fixed inverse frequency spectrum determined by a level-specific base scalar.
Specifically, for a head dimension $d$, the inverse frequency vector $\boldsymbol{\omega}^{(l)} \in \mathbb{R}^{d/2}$ in a given temporal level $l \in \{y, m, d\}$ is defined as
\begin{equation}
\boldsymbol{\omega}^{(l)}_i = \left(\text{base}^{(l)}\right)^{- \frac{2i}{d}}, \quad i = 0, \dots, \frac{d}{2}-1,
\end{equation}
where $\text{base}^{(l)}$ is a predefined scalar controlling the frequency range for temporal level $l$.
Given the decomposed temporal value $y_k, m_k, d_k$, the corresponding rotation angle is computed as
\begin{equation}
\boldsymbol{\theta}_k^{(y)} = y_k \cdot \boldsymbol{\omega}^{(y)}, \quad
\boldsymbol{\theta}_k^{(m)} = m_k \cdot \boldsymbol{\omega}^{(m)}, \quad
\boldsymbol{\theta}_k^{(d)} = d_k \cdot \boldsymbol{\omega}^{(d)}.
\end{equation}

The rotation angles are then used to inject temporal information into the query and key representations via rotary transformations.
Specifically, the query and key vectors are rotated in a pairwise manner over even-odd dimensions.
For each dimension pair $(2i, 2i+1)$, the rotation can be interpreted as applying a 2D rotation matrix parameterized by $\theta_{k,i}^{(l)}$:
\begin{equation}
\begin{pmatrix}
q_{k,2i}^{(l)} \\
q_{k,2i+1}^{(l)}
\end{pmatrix}
=
\begin{pmatrix}
\cos \theta_{k,i}^{(l)} & -\sin \theta_{k,i}^{(l)} \\
\sin \theta_{k,i}^{(l)} & \cos \theta_{k,i}^{(l)}
\end{pmatrix}
\begin{pmatrix}
q_{k,2i} \\
q_{k,2i+1}
\end{pmatrix},
\end{equation}
and the same transformation is applied to the key vector $\mathbf{k}_k$.

For notational simplicity, the rotary transformation can be written in a vectorized form as
\begin{equation}
\begin{aligned}
\mathbf{q}_k^{(l)} &= \mathbf{q}_k \odot \cos \boldsymbol{\theta}_k^{(l)} 
+ \mathrm{rotate}(\mathbf{q}_k) \odot \sin \boldsymbol{\theta}_k^{(l)}, \\
\mathbf{k}_k^{(l)} &= \mathbf{k}_k \odot \cos \boldsymbol{\theta}_k^{(l)} 
+ \mathrm{rotate}(\mathbf{k}_k) \odot \sin \boldsymbol{\theta}_k^{(l)},
\end{aligned}
\end{equation}
where $\mathrm{rotate}(\cdot)$ denotes a fixed permutation operator that rearranges the input vector by swapping each even-odd dimension pair.
Specifically, for an input vector $\mathbf{x} \in \mathbb{R}^{d}$, the operator is defined as
\begin{equation}
\mathrm{rotate}(\mathbf{x}) = (-x_1, x_0, -x_3, x_2, \dots, -x_{d-1}, x_{d-2}),
\end{equation}
which enables a vectorized implementation of the 2D rotation in each $(2i, 2i+1)$ subspace.This operation preserves vector norms while encoding temporal signals directly into the angular relationships between queries and keys, thereby enabling attention scores to be sensitive to relative temporal differences.

To capture temporal dynamics at multiple time scales, RoTE applies rotary transformations independently for the year-, month-, and day-level temporal components, producing three time-aware representations for each query and key.
These representations are then aggregated via a weighted fusion strategy:
\begin{equation}
\begin{aligned}
\mathbf{q}_k^{\text{time}} &=
\alpha_y \mathbf{q}_k^{(y)} +
\alpha_m \mathbf{q}_k^{(m)} +
\alpha_d \mathbf{q}_k^{(d)}, \\
\mathbf{k}_k^{\text{time}} &=
\alpha_y \mathbf{k}_k^{(y)} +
\alpha_m \mathbf{k}_k^{(m)} +
\alpha_d \mathbf{k}_k^{(d)},
\end{aligned}
\end{equation}
where $\alpha_y$, $\alpha_m$, and $\alpha_d$ are scalar weights that control the relative contributions of different temporal levels.
In our implementation, larger weights are assigned to coarse-grained temporal components to emphasize long-term preference evolution, while finer-grained components focus on short-term dynamics.



The resulting time-enhanced query and key representations are then used to compute attention scores using the standard scaled dot-product attention.
Importantly, RoTE does not modify the value representations, attention formulation, or training objective, and can be seamlessly integrated into existing Transformer-based sequential recommendation models.
This design makes RoTE a plug-and-play temporal modeling module that enhances temporal awareness without incurring significant computational overhead.


% \subsection{Integration with Sequential Recommendation Models}

% RoTE can be seamlessly integrated into both traditional and generative Transformer-based sequential recommendation models. 
% In practice, each item in the input sequence is augmented with its decomposed temporal features $(y_k, m_k, d_k)$, which serve as additional signals for temporal modeling. 
% Within the self-attention layers, the standard positional encoding is replaced by RoTE-based time-aware rotary transformations applied to the query and key representations. 
% By this way, RoTE can be applied broadly across different sequential recommendation frameworks, including both traditional paradigm and generative paradigm.

\section{Experiments}

\subsection{Experimental Setting}

\subsubsection{Dataset.}

We evaluate RoTE on three public benchmarks from the Amazon Reviews dataset \cite{amazon}: \textit{Sports and Outdoors}, \textit{Beauty}, and \textit{Toys and Games}. 
Following common practice in recent sequential recommendation research \cite{tiger}, we adopt the 5-core versions of these datasets, in which each user and each item has at least five interactions, thereby improves data quality and reduces noise in sparse interaction histories. 
Each dataset is chronologically sorted by timestamp, and sequences are truncated or padded to a fixed maximum length of 50 for both training and evaluation. Table \ref{tab:statistics} presents the detailed statistics of the three datasets.

\begin{table}[t]
    \centering
    \caption{Statistics of the Datasets.}
    \label{tab:statistics}
    \resizebox{\linewidth}{!}{
    \begin{tabular}{ccccc}
        \toprule
        Dataset & \#Users & \#Items & \#Interactions & \#Density \\
        \midrule
        Beauty & 22,363 & 12,101 & 198,360 & 0.00073 \\
        Sports and Outdoors & 35,598 & 18,357 & 296,175 & 0.00045 \\
        Toys and Games & 19,412 & 11,924 & 167,526 & 0.00073 \\
        \bottomrule
    \end{tabular}
    }
\end{table}


\subsubsection{Baselines.}

We compare RoTE with the following baseline models, covering both traditional and generative paradigms. 
Traditional baselines \cite{gru4rec, caser, bert4rec, sasrec} capture sequential dependencies in user behavior via next-item prediction, while generative baselines \cite{vqrec, tiger, hstu, rpg} adopt autoregressive generation frameworks to model user interaction sequences. 
To evaluate the effectiveness of RoTE as a plug-and-play temporal modeling module, we integrate it into selected representative Transformer-based backbones from each paradigm, while other baselines are reported in their original forms. 

\begin{table*}[t]
\caption{Overall performance comparison of state-of-the-art sequential recommendation models and their RoTE-enhanced counterparts on three benchmark datasets. RoTE is a plug-and-play module that improves model performance without modifying the backbone architectures. Performance improvements of RoTE-based models over their corresponding backbones are statistically significant according to a paired t-test with $p < 0.05$.}
\centering
\setlength{\tabcolsep}{4pt}
\begin{tabular}{
  l|l % 这条 | 会应用到表格的每一行
  c c c c
  c c c c
  c c c c
}
\toprule
% --- 表头 ---
\multicolumn{1}{l}{} & \multirow{2.5}{*}{\textbf{Methods}} & \multicolumn{4}{c}{\textbf{Sports and Outdoors}} & \multicolumn{4}{c}{\textbf{Beauty}}  & \multicolumn{4}{c}{\textbf{Toys and Games}} \\
\cmidrule(lr){3-6}\cmidrule(lr){7-10}\cmidrule(lr){11-14}
\multicolumn{1}{l}{} & & R@5 & N@5 & R@10 & N@10 & R@5 & N@5 & R@10 & N@10 & R@5 & N@5 & R@10 & N@10 \\
\midrule
% --- Traditional 组 ---
\multirow{6}{*}{\rotatebox{90}{\textbf{Traditional}}}
& GRU4Rec     & 0.0129 & 0.0086 & 0.0204 & 0.0110 & 0.0164 & 0.0099 & 0.0283 & 0.0137 & 0.0097 & 0.0059 & 0.0176 & 0.0084 \\
& Caser       & 0.0116 & 0.0072 & 0.0194 & 0.0097 & 0.0205 & 0.0131 & 0.0347 & 0.0176 & 0.0166 & 0.0107 & 0.0270 & 0.0141 \\
& BERT4Rec    & 0.0115 & 0.0075 & 0.0191 & 0.0099 & 0.0203 & 0.0124 & 0.0347 & 0.0170 & 0.0116 & 0.0071 & 0.0203 & 0.0099 \\
& SASRec      & 0.0228 & 0.0146 & 0.0347 & 0.0189 & 0.0375 & 0.0238 & 0.0604 & 0.0311 & 0.0466 & 0.0310 & 0.0688 & 0.0375 \\
& \textbf{RoTE-SASRec} & \textbf{0.0233} & \textbf{0.0158} & \textbf{0.0372} & \textbf{0.0202} & \textbf{0.0400} & \textbf{0.0260} & \textbf{0.0638} & \textbf{0.0336} & \textbf{0.0512} & \textbf{0.0345} & \textbf{0.0747} & \textbf{0.0420} \\
& \textit{Improv.} & \textcolor{green!50!black}{\textit{+2.19\%}} & \textcolor{green!50!black}{\textit{+8.22\%}} & \textcolor{green!50!black}{\textit{+7.20\%}} & \textcolor{green!50!black}{\textit{+6.88\%}} & \textcolor{green!50!black}{\textit{+6.67\%}} & \textcolor{green!50!black}{\textit{+9.24\%}} & \textcolor{green!50!black}{\textit{+5.63\%}} & \textcolor{green!50!black}{\textit{+8.04\%}} & \textcolor{green!50!black}{\textit{+9.87\%}} & \textcolor{green!50!black}{\textit{+11.29\%}} & \textcolor{green!50!black}{\textit{+8.58\%}} & \textcolor{green!50!black}{\textit{+12.00\%}} \\
\midrule
% --- Generative 组 ---
\multirow{6}{*}{\rotatebox{90}{\textbf{Generative}}}
& VQ-Rec      & 0.0208 & 0.0137 & 0.0300 & 0.0173 & 0.0457 & 0.0317 & 0.0664 & 0.0383 & 0.0373 & 0.0256 & 0.0484 & 0.0323 \\
& TIGER       & 0.0264 & 0.0181 & 0.0400 & 0.0225 & 0.0454 & 0.0321 & 0.0648 & 0.0384 & 0.0521 & 0.0371 & 0.0712 & 0.0432 \\
& HSTU        & 0.0258 & 0.0165 & 0.0414 & 0.0215 & 0.0469 & 0.0314 & 0.0704 & 0.0389 & 0.0433 & 0.0281 & 0.0669 & 0.0357 \\
& RPG         & 0.0294 & 0.0201 & 0.0419 & 0.0241 & 0.0504 & 0.0349 & 0.0730 & 0.0422 & 0.0531 & 0.0373 & 0.0759 & 0.0446 \\
& \textbf{RoTE-RPG}    & \textbf{0.0317} & \textbf{0.0222} & \textbf{0.0456} & \textbf{0.0266} & \textbf{0.0521} & \textbf{0.0370} & \textbf{0.0742} & \textbf{0.0438} & \textbf{0.0624} & \textbf{0.0448} & \textbf{0.0852} & \textbf{0.0521} \\
& \textit{Improv.} & \textcolor{green!50!black}{\textit{+7.82\%}} & \textcolor{green!50!black}{\textit{+10.45\%}} & \textcolor{green!50!black}{\textit{+8.83\%}} & \textcolor{green!50!black}{\textit{+10.37\%}} & \textcolor{green!50!black}{\textit{+3.37\%}} & \textcolor{green!50!black}{\textit{+6.02\%}} & \textcolor{green!50!black}{\textit{+1.64\%}} & \textcolor{green!50!black}{\textit{+3.79\%}} & \textcolor{green!50!black}{\textit{+17.51\%}} & \textcolor{green!50!black}{\textit{+20.11\%}} & \textcolor{green!50!black}{\textit{+12.26\%}} & \textcolor{green!50!black}{\textit{+16.82\%}} \\
\bottomrule
\end{tabular}
\label{tab:main_results}
\end{table*}

\subsubsection{Implementation Details.}

We implement RoTE module using the PyTorch framework. 
For generative models, we adopt RPG \cite{rpg} as a representative backbone and build our implementation upon its official codebase\footnote{https://github.com/facebookresearch/RPG\_KDD2025}. 
For traditional models, we select SASRec \cite{sasrec, sasrec_revis} as a representative and use the publicly available implementation\footnote{https://github.com/pmixer/SASRec.pytorch}. Based on these backbones, we replace the original position encoding module with RoTE while keeping the remaining architecture unchanged. 
For the year-, month-, and day-level temporal components, the base frequencies are set to $\text{base}^{(y)} = 10^6$, $\text{base}^{(m)} = 10^4$, and $\text{base}^{(d)} = 10^2$, respectively, and the corresponding fusion weights are set to $\alpha_y = 1.5$, $\alpha_m = 1.0$, and $\alpha_d = 0.5$. 
Except for the integration of RoTE, all other architectural configurations follow the standard setups reported in the original implementations.

\subsubsection{Evaluation Metrics.}

We evaluate model performance using Recall@K and NDCG@K, where $K \in \{5, 10\}$. 
Following the standard leave-one-out evaluation strategy, we reserve the last interaction of each user as the test item, the penultimate item for validation, and all remaining interactions for training. Results are averaged over five random seeds.

\subsection{Overall Performance}


Table \ref{tab:main_results} reports the overall performance of representative sequential recommendation models and their RoTE-enhanced variants on three Amazon benchmarks.
As shown in the results, incorporating RoTE consistently improves the performance of both traditional and generative models across all datasets and evaluation metrics.
Specifically, RoTE brings stable gains over strong baselines, including SASRec and RPG, achieving noticeable improvements in performance. Notably, on the \textit{Toys and Games} benchmark, RoTE improves Recall@5 and NDCG@5 by 17.51\% and 20.11\% over the RPG baseline, respectively.
These improvements indicate that RoTE effectively captures temporal information and can be seamlessly integrated into diverse sequential modeling paradigms.
Overall, the results demonstrate the effectiveness and broad applicability of RoTE as a plug-and-play module for sequential recommendation.

\subsection{Ablation Study}

\begin{table}[t]
\centering
\caption{Ablation study of different levels of temporal signals on \textit{Toys and Games} dataset.}
\label{tab:ablation}
\begin{tabular}{lccccc}
\toprule
\multirow{2}{*}{\textbf{Methods}} 
& \multicolumn{2}{c}{\textbf{SASRec}} 
& \multicolumn{2}{c}{\textbf{RPG}} \\
\cmidrule(lr){2-3} \cmidrule(lr){4-5}
& \textbf{R@10} & \textbf{N@10} 
& \textbf{R@10} & \textbf{N@10} \\
\midrule
Backbone      & 0.0688 & 0.0375 & 0.0759 & 0.0446 \\
+ Pure Timestamp          & 0.0713 & 0.0403 & 0.0838 & 0.0498 \\
+ Y                       & 0.0701 & 0.0387 & 0.0815 & 0.0467 \\
+ Y + M                     & 0.0725 & 0.0409 & 0.0846 & 0.0505 \\
\textbf{+ Y + M + D (RoTE)}      & \textbf{0.0747} & \textbf{0.0420} & \textbf{0.0852} & \textbf{0.0521} \\
\bottomrule
\end{tabular}
\end{table}

Table \ref{tab:ablation} reports ablation results with different temporal encoding strategies on \textit{Toys and Games}. 
\textbf{Backbone} denotes the original models using positional embeddings, which only preserve relative order and can be viewed as a degenerate form of temporal modeling. 
\textbf{Pure Timestamp} extends this formulation by directly applying rotary embedding to absolute Unix timestamps.
The \textbf{Y}, \textbf{Y+M}, and \textbf{Y+M+D} variants progressively introduce year-, month-, and day-level calendar components into RoTE. 
Results show that while absolute timestamps provide measurable improvements over positional encoding, performance consistently improves as finer-grained and semantically structured temporal components are incorporated. 
\textbf{Y+M+D} achieves the best results on both SASRec and RPG, validating the effectiveness of coarse-to-fine multi-level temporal modeling.

\subsection{Efficiency Analysis}

\begin{table}[t]
\centering
\caption{Efficiency analysis of RoTE on different backbones.}
\label{tab:efficiency}
\begin{tabular}{lccc}
\toprule
\textbf{Methods} & \textbf{\#Params} & \textbf{FLOPs} & \textbf{Latency} \\
\midrule
SASRec            & 817.02K & 6.21M & 1.686ms \\
RoTE-SASRec     & 813.76K & 6.32M & 2.423ms \\
\midrule
RPG               & 13.58M & 994.94M & 6.224ms \\
RoTE-RPG         & 13.55M & 994.97M & 8.154ms \\
\bottomrule
\end{tabular}
\end{table}

Table \ref{tab:efficiency} compares the efficiency of backbone models and their RoTE-enhanced variants. RoTE introduces negligible computational overhead, with only marginal increases in FLOPs and inference latency, while maintaining millisecond-level real-time performance. Notably, the slight reduction in parameter count stems from replacing the original positional embedding table with a lightweight computation-based temporal encoding that does not rely on large embedding lookups. Overall, RoTE enables Transformer-based sequential recommendation models to better capture heterogeneous temporal distances among user interactions without incurring significant additional computational cost.

\section{Conclusion}

In this work, we propose RoTE, a plug-and-play rotary time embedding module for Transformer-based sequential recommendation. RoTE models time in a hierarchical and semantically structured manner, enabling models to better capture heterogeneous temporal distances among user interactions. Extensive experiments demonstrate that RoTE consistently improves performance across both traditional and generative backbones while introducing negligible computational overhead. These results indicate that structured temporal modeling is a practical and effective direction for enhancing sequential recommender systems.

\section*{Acknowledgements}

This work was supported by the STI 2030-Major Projects under Grant 2021ZD0201404  and Shenzhen Key Laboratory of New Generation Interactive Media Technology Innovation (ZDSYS2021062\\3092001004).

% \clearpage

%%
%% The next two lines define the bibliography style to be used, and
%% the bibliography file.
\bibliographystyle{ACM-Reference-Format}
\bibliography{SIGIR2026/references}


%%
%% If your work has an appendix, this is the place to put it.


\end{document}
\endinput
%%
%% End of file `sample-sigconf.tex'.
